\subsection{POLYNOMIAL ALGEBRAS}

Let B be a unital commutative associative A-algebra and let \((x_{i})_{i\in I}\) be a family of elements of B; the subalgebra of B generated by the \(x_{i}\) ( \(i\in I\) ) and the unit element is denoted by \(\mathrm{A}[(x_{i})_{i\in I}]_{\mathrm{B}}\) or simply \(\mathrm{A}[(x_{i})_{i\in I}]\) when no confusion can arise. For every set I, the algebra \(\operatorname{Libasc}_{\mathbf{A}}(\mathbf{I})\) is therefore equal to \(\mathrm{A}[(\mathbf{X}_{i})_{i\in I}]\) (also denoted by \(\mathrm{A}[\mathbf{X}_{i}]_{i\in I}\) ); the latter notation, which has the advantage of indicating the notation chosen to denote the indeterminates, is the one we shall generally use in the rest of this Treatise. The elements of \(\mathrm{A}[(\mathbf{X}_{i})_{i\in I}]\) are called polynomials in the indeterminates \(X_{i}\) ( \(i\in I\) ) with coefficients in A; it is a convention that, when it is said ``let \(\mathrm{A}[(\mathbf{X}_{i})_{i\in I}]\) be a polynomial algebra'', the \(X_{i}\) are always understood to be the indeterminates. For every subset J of I, the use of the above notation amounts to identifying \(\operatorname{Libasc}_{\mathbf{A}}(\mathbf{J})\) with the algebra of \(\operatorname{Libasc}_{\mathbf{A}}(\mathbf{I})\) generated by the \(X_{i}\) of index \(i\in J\) and the unit element (cf.~no. 7, Remark (3)). For \(I=\{1,2,\ldots,n\}\) , we write \(\mathrm{A}[\mathbf{X}_{1},\mathbf{X}_{2},\ldots,\mathbf{X}_{n}]\) instead of \(\mathrm{A}[(\mathbf{X}_{i})_{i\in I}]\) .

If I and I' are two equipotent sets, the algebras \(\operatorname{Libasc}_{\mathbf{A}}(\mathbf{I})\) and \(\operatorname{Libasc}_{\mathbf{A}}(\mathbf{I}')\) are isomorphic. A{[}X{]} is often used to denote the polynomial algebra corresponding to an unspecified indexing set with a single element, X denoting the unique indeterminate; similarly, A{[}X, Y{]}, A{[}X, Y, Z{]}, \ldots{} are used to denote the polynomial algebras corresponding to unspecified indexing sets with

2, 3, \ldots{} elements. Note that, by virtue of the conventions made above, A{[}X{]} and A{[}Y{]} are for example (distinct) subalgebras of A{[}X, Y, Z{]} if A ≠ \{0\}.

The elements

\[
\mathrm{X} ^ {\nu} = \prod_ {i \in I} \mathrm{X} _ {i} ^ {\nu (i)},
\]

where v runs through \(\mathbf{N}^{(1)}\) form a basis of the polynomial algebra \(\mathrm{A}[(\mathrm{X}_{i})_{i\in\mathbb{I}}]\) . These elements are called monomials in the indeterminates \(X_{i}\) and the number \(|v|=\sum_{i\in\mathbb{I}}v(i)\) is called the degree (or total degree) of the monomial \(X^{v}\) . The unique monomial of degree 0 is the unit element of \(\mathrm{A}[(\mathrm{X}_{i})_{i\in\mathbb{I}}]\) ; it is often identified with the unit element 1 of A. Every polynomial u of \(\mathrm{A}[(\mathrm{X}_{i})_{i\in\mathbb{I}}]\) can be written uniquely as

\[
u = \sum_ {v \in \mathbf {N} ^ {(I)}} \alpha_ {v} X ^ {v}
\]

with \(\alpha_{\nu} \in \mathbf{A}\) ; the elements \(\alpha_{\nu}\) , zero except for a finite number of indices \(\nu \in \mathbf{N}^{(1)}\) , are called the coefficients of \(u\) ; the elements \(\alpha_{\nu} \mathrm{X}^{\nu}\) are called the terms of \(u\) (the element \(\alpha_{\nu} \mathrm{X}^{\nu}\) often being called the ``term in \(\mathrm{X}^{\nu''}\) ); in particular the term \(\alpha_{0} \mathrm{X}^{0}\) (identified with \(\alpha_{0} \in \mathbf{A}\) ) is called the constant term of \(u\) . If \(J\) is a subset of \(I\) , \(u\) belongs to \(\mathrm{A}[(\mathrm{X}_{i})_{i \in J}]\) if and only if \(\alpha_{\nu} = 0\) for \(\nu \notin \mathbf{N}^{(J)}\) . It follows that \(\mathrm{A}[(\mathrm{X}_{i})_{i \in I}]\) is the union of the subalgebras \(\mathrm{A}[(\mathrm{X}_{i})_{i \in J}]\) , where \(J\) runs through the set of finite subsets of \(I\) . If \(\alpha_{\nu} = 0\) for \(|\nu| > n\) , \(u\) is said to be a polynomial of degree \(\leqslant n\) . When \(\alpha_{\nu} = 0\) , (by an abuse of language) \(u\) is said to contain no term in \(\mathrm{X}^{\nu'}\) ; in particular, when \(\alpha_{0} = 0\) , \(u\) is said to be a polynomial with no constant term.

For every non-zero polynomial \(u = \sum_{v} \alpha_v X^v\) , the degree (or total degree) of \(u\) is the greatest of the integers \(|v|\) of the multiindices \(v\) such that \(\alpha_v \neq 0\) .

Let F be a unital associative A-algebra and let \((x_{i})_{i\in I}\) be a family of elements of F, which are pairwise permutable. The subalgebra \(F'\) of F generated by the \(x_{i}\) and the unit element is commutative (§1, no. 7), which allows us to define substituting the \(x_{i}\) for the \(X_{i}\) in the polynomial \(u\in\mathrm{A}[(X_{i})_{i\in I}]\) (although F is not necessarily commutative): \(u((x_{i})_{i\in I})\) is an element of \(F'\) and therefore of F and \(h\colon u\to u((x_{i})_{i\in I})\) is a homomorphism of \(\mathrm{A}[(X_{i})_{i\in I}]\) into F. The elements of the kernel of h are the relators of the family \((x_{i})\) in \(\mathrm{A}[(X_{i})_{i\in I}]\) , also called polynomial relators (with coefficients in A) between the \(x_{i}\) . The image of the homomorphism h is the subalgebra \(F'\) , also denoted by \(\mathrm{A}[(x_{i})_{i\in I}]\) (even when F is not commutative); if a is the ideal of polynomial relators between the \(x_{i}\) , then there is an exact sequence of A-modules

\[
0 \longrightarrow a \longrightarrow A [ (X _ {i}) _ {i \in I} ] \stackrel {h} {\longrightarrow} A [ (x _ {i}) _ {i \in I} ] \longrightarrow 0.
\]

PROPOSITION 8. Let \(\mathbf{A}[(\mathbf{X}_i)_{i\in I}]\) be a polynomial algebra, \(J\) a subset of \(I\) and \(K\) the complement of \(J\) in \(I\) . Writing \(\mathbf{A}' = \mathbf{A}[(\mathbf{X}_j)_{j\in J}]\) and denoting by \(\mathbf{X}_k'\) ( \(k\in K\) ) the indeterminates in the polynomial algebra \(\operatorname{Libasc}_{\mathbf{A}'}(\mathbf{K}) = \mathbf{A}'[(\mathbf{X}_k')_{k \in \mathbb{K}}]\) , there exists a unique ring isomorphism of \(\mathbf{A}'[(\mathbf{X}_k')_{k \in \mathbb{K}}]\) onto \(\mathbf{A}[(\mathbf{X}_i)_{i \in I}]\) which coincides with the identity on \(\mathbf{A}'\) and maps \(\mathbf{X}_k'\) to \(\mathbf{X}_k\) for all \(k \in \mathbb{K}\) .

Clearly \(\mathbf{A}[(\mathbf{X}_i)_{i\in I}]\) is an \(\mathbf{A}'\) -algebra generated by the \(\mathbf{X}_k\) for \(k\in \mathbf{K}\) . On the other hand, as a polynomial relator between the \(\mathbf{X}_k\) ( \(k\in \mathbf{K}\) ) with coefficients in \(\mathbf{A}'\) can be written uniquely as \(\sum_{\nu}h_{\nu}((\mathbf{X}_j)_{j\in J})\mathbf{X}^{\nu}\) where \(\nu\) runs through a finite subset of \(\mathbf{N}^{(\mathbb{R})}\) and where the \(h_\nu\) are elements of \(\mathbf{A}[(\mathbf{X}_j)_{j\in J}]\) , the \(h_\nu\) must be polynomial relators between the \(\mathbf{X}_j\) with coefficients in \(\mathbf{A}\) and hence are all zero, which proves the proposition.

The isomorphism described in Proposition 8 is often used to identify the elements of \(\mathbf{A}[(\mathbf{X}_i)_{i\in I}]\) with polynomials with coefficients in \(\mathbf{A}' = \mathbf{A}[(\mathbf{X}_j)_{j\in J}]\) . If \(u\) is an element \(\neq 0\) of \(\mathbf{A}[(\mathbf{X}_i)_{i\in I}]\) , its total degree considered as an element of \(\mathbf{A}'[(\mathbf{X}_k)_{k\in K}]\) is also called its degree with respect to the \(\mathbf{X}_i\) of index \(i\in K\) .

Remark. Let I and J be two sets and \((\mathrm{P}_{j})_{j\in J}\) a family of elements of \(\mathbf{Z}[(\mathrm{X}_{i})_{i\in I}]\) ; if Q is an element of \(\mathbf{Z}[(\mathrm{X}_{j})_{j\in J}]\) such that \(\mathbf{Q}((\mathrm{P}_{j})_{j\in J})=0\) , then, for every family \((b_{i})_{i\in I}\) of pairwise permutable elements of a ring B,

\[
\mathrm{Q} ((\mathrm{P} _ {j} ((b _ {i}) _ {i \in \mathrm{I}})) _ {j \in \mathrm{J}}) = 0.
\]

The relations which hold of the form \(\mathbf{Q}((\mathbf{P}_{j})_{j\in\mathbf{J}})=0\) are sometimes called polynomial identities. For example

\[
\left(\mathrm{X} _ {1} + \mathrm{X} _ {2}\right) ^ {2} - \mathrm{X} _ {1} ^ {2} - \mathrm{X} _ {2} ^ {2} - 2 \mathrm{X} _ {1} \mathrm{X} _ {2} = 0
\]

with

\[
\begin{array}{r l} \mathrm{Q} & = \mathrm{Y} _ {1} ^ {2} - \mathrm{Y} _ {2}, \qquad \mathrm{P} _ {1} = \mathrm{X} _ {1} + \mathrm{X} _ {2}, \qquad \mathrm{P} _ {2} = \mathrm{X} _ {1} ^ {2} + \mathrm{X} _ {2} ^ {2} + 2 \mathrm{X} _ {1} \mathrm{X} _ {2} \\ & \quad \mathrm{X} _ {1} ^ {n} - \mathrm{X} _ {2} ^ {n} - (\mathrm{X} _ {1} - \mathrm{X} _ {2}) (\mathrm{X} _ {1} ^ {n - 1} + \mathrm{X} _ {1} ^ {n - 2} \mathrm{X} _ {2} + \dots + \mathrm{X} _ {2} ^ {n - 1}) = 0 \end{array}
\]

with

\[
\begin{array}{c} \mathrm{Q} = \mathrm{Y} _ {1} - \mathrm{Y} _ {2} \mathrm{Y} _ {3}, \qquad \mathrm{P} _ {1} = \mathrm{X} _ {1} ^ {n} - \mathrm{X} _ {2} ^ {n}, \qquad \mathrm{P} _ {2} = \mathrm{X} _ {1} - \mathrm{X} _ {2}, \\ \mathrm{P} _ {3} = \mathrm{X} _ {1} ^ {n - 1} + \mathrm{X} _ {1} ^ {n - 2} \mathrm{X} _ {2} + \dots + \mathrm{X} _ {2} ^ {n - 1} \end{array}
\]

are polynomial identities.

